\section{Round once, with an explicit tie rule}
Add 16384 to the complete accumulator, then shift right arithmetically by fifteen
with \texttt{dsra}. Arithmetic right shift gives floor division by a power of two.
The bias is half the divisor, so the combination implements the specified nearest
rounding rule for both signs. The bias is also safe within the established 64-bit
bound. A logical shift would turn a negative value into a large positive integer.

\begin{table}[htbp]
\centering
\caption{Rounding boundaries in accumulator units, before output saturation.}
\tagpdfsetup{table/header-rows={1}}
\begin{tabular}{rrr}
\toprule
\(S\) & \(S+16384\) & Rounded sample code\\
\midrule
-49153 & -32769 & -2\\
-49152 & -32768 & -1\\
-16385 & -1 & -1\\
-16384 & 0 & 0\\
16383 & 32767 & 0\\
16384 & 32768 & 1\\
49152 & 65536 & 2\\
\bottomrule
\end{tabular}
\end{table}

Halfway values go toward positive infinity: negative one half becomes zero, and
negative one and a half becomes negative one. This is not nearest-even rounding or
rounding every tie away from zero. Ordinary signed integer division that truncates
toward zero also differs from the required floor operation for negative biased values.

Do not round each product before adding. With eight samples equal to one and eight
coefficients equal to 16384, each scaled product is one half. Rounding each separately
would produce eight. Accumulating the exact products first gives four. A fixed-point
multiply instruction that scales or saturates each lane early therefore needs a separate
proof; its name alone does not make it suitable for this contract.
